% ===================================================================
%  TABELL Nr 8 — Skörden. — Jakten, vinskörden, fisket.
%  Traduction 2026 d'apres texto/io, controlee sur texto/fr et
%  texto/es. Consigne : voir texto/sv/10-tabell-01.tex.
%
%  L'appel de note est dans le TITRE de la scene. LE RENVOI (13) DE
%  LA FAUX MANQUE A L'IDO ; la colonne le rend : ecart declare.
% ===================================================================

%%K t08-noto-1 noto
\VUnotes{143.2mm}{%
\textit{(*) Eftersom det ordet inte är noggrant : rör det sig om linskörd,
vinskörd, äppelskörd? Kanske vore det bra att använda «} \VUgras{mesono}
\textit{» som föreslogs i Mondo XIV, s. 242. Men hittills har den nya roten inte
antagits av Akademien och väntar på diskussioner enligt det vanliga idistiska
reglementet.}%
}

%%K t08-apar-1 apar
\VUsaut{5.0mm}
\VUtitre{15.0pt}{60}{TABELL Nr 8}
\VUsaut{3.0mm}
\VUcentre{13.2pt}{90}{\textit{Första scenen.}}
\VUsaut{3.0mm}
\VUpk{10.2pt}{40}{Skörden (*).}
\VUsaut{4.0mm}
\VUpk{10.2pt}{40}{Landsbygdens utseenden.}
\VUsaut{3.0mm}

%%K t08-c1-01-1 p
1. --- Med \VUgras{sommaren} har landsbygdens utseende ännu en gång förändrats.
Utsädet som anförtroddes åt fåran har grott; en grön planta kom upp ur marken och
gick i ax. Kornet bildades, sedan mognade det och nu behöver man bara skörda.

%%K t08-c1-02-1 p
2. --- \VUgras{Ängsblommorna} har slagit ut. \VUgras{Blåklint}
\textsuperscript{(1)}, \VUgras{vallmo} \textsuperscript{(2)} och
\VUgras{fingerborgsblomma} \textsuperscript{(3)} blandar i vetefältens enformiga
färgton den glada kontrasten av sina friska \VUgras{blomkronor}.

%%K t08-c1-03-1 p
3. --- Fruktträden har smyckat sig med löv; frukten har mognat. Det lilla
\VUgras{körsbärsträdet} \textsuperscript{(4)} är fullt av vackra
\VUgras{körsbär} \textsuperscript{(5)} som barnen med stort nöje plockar och
äter.

%%K t08-c1-04-1 p
4. --- Nu kan boskapen tillbringa hela dagen i fria luften.

%%K t08-c1-04-2 p
\VUgras{Lammen} \textsuperscript{(6)} har vuxit och blivit vackra
\VUgras{tackor} \textsuperscript{(7)} och vackra \VUgras{baggar}
\textsuperscript{(8)} med tät ull. De betar lugnt \VUgras{gräset} på
\VUgras{betesmarken} \textsuperscript{(9)} som är prickig av
\VUgras{tusenskönor}.

%%K t08-c1-04-3 p
Snart skall man börja klippa de djuren för att få deras \VUgras{ull} av vilken
man gör klädet till våra kläder.

%%K t08-tit-2 sub
\VUsaut{5.0mm}
\VUpk{10.2pt}{40}{Herden.}
\VUsaut{3.4mm}

%%K t08-c1-05-1 p
5. --- \VUgras{Herden} \textsuperscript{(10)} tycks inte ge dem särskilt mycket
akt. Han bryr sig bara om ovädret som drar förbi vid horisonten, och
\VUgras{fårvaktaren} \textsuperscript{(13)}, som för en \VUgras{fältflaska}
\textsuperscript{(14)} till sina läppar verkar ännu mer sorglös.

%%K t08-c1-06-1 p
6. --- Hettan är tryckande; ty \VUgras{hunden} \textsuperscript{(15)} kommer,
också han, för att släcka sin törst; han lapar det friska vattnet i
\VUgras{bäcken} \textsuperscript{(16)} och skrämmer en stackars \VUgras{liten
groda} \textsuperscript{(17)} som hade gömt sig i strandgräset. Denna hoppar i
det klara vattnet där \VUgras{näckrosorna} \textsuperscript{(18)} blommar.

%%K t08-c1-07-1 p
7. --- En \VUgras{ärla} \textsuperscript{(19)} badar vid den lilla
\VUgras{källan} \textsuperscript{(20)} som sprutar fram mellan två klippor och
gömmer sig under \VUgras{törnbuskarna} \textsuperscript{(21)} och
\VUgras{hagtornsbuskarna} \textsuperscript{(22)}.

%%K t08-tit-3 sub
\VUsaut{5.0mm}
\VUpk{10.2pt}{40}{Skörden.}
\VUsaut{3.4mm}

%%K t08-c1-08-1 p
8. --- För lantbarnen är sädesskörden en mycket rolig tid. De gör glatt
\VUgras{kullerbyttor} \textsuperscript{(24)} på \VUgras{halmhögarna}
\textsuperscript{(23)}, de hjälper \VUgras{skördemännen} \textsuperscript{(26)}
och medan dessa slår \VUgras{sädesfältet} \textsuperscript{(27)} med sina
\VUgras{liar} \textsuperscript{(25)} som är väl vässade med \VUgras{brynstenen}
\textsuperscript{(30)}, samlar de säden och binder den till \VUgras{kärvar}
\textsuperscript{(31)} eller till \VUgras{små knippen} \textsuperscript{(32)}.

%%K t08-c1-09-1 p
9. --- Ibland tillåter man de största bland dem att skära av med en
\VUgras{skära} \textsuperscript{(12)} de \VUgras{gyllene stråna}
\textsuperscript{(28)} som bär de tunga \VUgras{axen} \textsuperscript{(29)}
fulla av korn; och de små, som vaktar \VUgras{boskapen} \textsuperscript{(33)}
som betar på \VUgras{betesmarken} \textsuperscript{(34)} under skuggan av den
stora \VUgras{linden} \textsuperscript{(35)} fångar med sin \VUgras{håv}
\textsuperscript{(36)} de vackra \VUgras{fjärilarna}.

%%K t08-c1-09-2 p
Några klättrar till och med upp på \VUgras{kärran} \textsuperscript{(37)} för att
ordna kärvarna som man räcker dem med en \VUgras{gaffel}
\textsuperscript{(38)}.

%%K t08-c1-10-1 p
10. --- På hela \VUgras{slätten} \textsuperscript{(39)} där \VUgras{lärkorna}
\textsuperscript{(41)} börjar flyga upp råder livligheten. Alla är sysselsatta
med skörden. Men medan de fattiga fortsätter att använda de gamla redskapen,
\VUgras{liar, skäror} och \VUgras{slagor} \textsuperscript{(40)}, låter de rika
godsägarna slå sitt vete eller sin \VUgras{råg} \textsuperscript{(43)} med en
\VUgras{slåttermaskin} \textsuperscript{(42)}. Sedan gör man, eftersom man inte
behöver bära kärvarna till ladan, stora \VUgras{kärvhögar}
\textsuperscript{(44)}.

%%K t08-c1-11-1 p
11. --- Då för man dit en \VUgras{tröskmaskin} \textsuperscript{(47)} som ett
\VUgras{lokomobil} \textsuperscript{(45)} driver med en \VUgras{drivrem}
\textsuperscript{(46)} och så skiljs kornet från \VUgras{halmen}, utan att lämna
åkern där det såddes.

%%K t08-tit-4 sub
\VUsaut{5.0mm}
\VUpk{10.2pt}{40}{Ovädret.}
\VUsaut{3.4mm}

%%K t08-c1-12-1 p
12. --- Ofta kommer svåra \VUgras{oväder} \textsuperscript{(53)}, under
skördetiden. Först hörs \VUgras{åskans} mullrande på avstånd; en \VUgras{häftig
västanvind} \textsuperscript{(54)} börjar blåsa och böjer flåsande de grövsta träden i
\VUgras{skogen} \textsuperscript{(49)}. Svarta \VUgras{moln}
\textsuperscript{(56)} stormar himlen; de genomkorsas av \VUgras{blixtens}
\textsuperscript{(55)} bländande sicksacklinjer. \VUgras{Regnet} och ibland
\VUgras{hagel} börjar falla, och slår ner skörden som är färdig att skäras.

%%K t08-c1-13-1 p
13. --- Ofta antänder blixten en byggnad. Så lades förra året \VUgras{väderkvarnens}
\textsuperscript{(50)} tak i aska och två av dess \VUgras{vingar}
\textsuperscript{(51)} krossades.

%%K t08-c1-14-1 p
14. --- Efter några timmar återvänder lugnet. En strålande \VUgras{regnbåge}
\textsuperscript{(52)} visar sig vid horisonten och naturen återtar sitt liv som
avbrutits några ögonblick. Lantborna, som hade flytt in i \VUgras{hyddorna}
\textsuperscript{(48)} börjar arbeta igen och söker modigt reparera de skador de
just har lidit.

%%K t08-tit-5 sub
\VUsaut{6.0mm}
\VUcentre{13.2pt}{90}{\textit{Andra scenen.}}
\VUsaut{3.6mm}

%%K t08-tit-6 sub
\VUpk{10.2pt}{40}{Jakten. --- Vinskörden.}
\VUsaut{1.4mm}
\VUpk{10.2pt}{40}{Fisket.}
\VUsaut{4.0mm}

%%K t08-c2-01-1 p
1. --- Omkring slutet av \VUgras{augusti} eller mitten av \VUgras{september},
allt efter trakterna, äger jaktens öppnande rum. Knappt har de första
solstrålarna fått \VUgras{ödlorna} \textsuperscript{(2)} att komma ut ur sitt
hål, förrän \VUgras{jägaren} \textsuperscript{(5)} ger sig i väg med sina
\VUgras{hundar} \textsuperscript{(4)}.

%%K t08-c2-02-1 p
2. --- Han har tagit på sig sin stadiga \VUgras{jaktdräkt}
\textsuperscript{(9)} av tjockt kläde, satt på sig \VUgras{damasker} av läder,
med vilka han skall lyckas gå i det fuktiga gräset där \VUgras{paddorna}
\textsuperscript{(1)} gömmer sig; han har tagit sitt \VUgras{gevär}
\textsuperscript{(6)} och sin \VUgras{jaktväska} \textsuperscript{(10)}, och här
är det att han far i väg.

%%K t08-c2-03-1 p
3. --- Förföljelsens iver för honom mitt in i en vingård där vinskörden är mycket
livlig. Vinodlarens barn, som vanligen vaktar getterna vid vägen, låter dem i dag
beta på måfå, och sedan de har bundit en gammal \VUgras{bock}
\textsuperscript{(3)} vid en påle som inte skulle underlåta att begagna sin
frihet för att rymma, tillägnar de sig i sin tur en del i nöjet. Den ene griper
en druvklase i en \VUgras{skördekorg} \textsuperscript{(12)} och roar sig med
att låta \VUgras{saften} \textsuperscript{(14)} rinna i en liten bägare
\textsuperscript{(13)} för att göra must (ojäst vin). En annan pryder sig med en
\VUgras{lövkrans} \textsuperscript{(15)} som han just har flätat. Bredvid dem
kommer oförskämda \VUgras{sparvar} \textsuperscript{(16)} ända fram till pressen
för att röva druvorna.

%%K t08-c2-04-1 p
4. --- Medan barnen roar sig, arbetar männen. \VUgras{Vinskördarna}
\textsuperscript{(18)} bär druvorna in i källaren med \VUgras{ryggkorgar}
\textsuperscript{(17)}; en \VUgras{tunnbindare} \textsuperscript{(19)} lagar
\VUgras{tunnorna} \textsuperscript{(20)}, kontrollerar om \VUgras{laggarna}
\textsuperscript{(22)} och \VUgras{bottnarna} \textsuperscript{(21)} är i gott
skick, och byter ut de brustna \VUgras{banden} \textsuperscript{(23)}. Han har
också arbetat på \VUgras{karen} \textsuperscript{(25)} i vilka man häller
\VUgras{druvorna} \textsuperscript{(26)} för att låta dem jäsa.

%%K t08-c2-05-1 p
5. --- När jäsningen har slutat, tappar man vinet och lägger återstoden på
\VUgras{pressen} \textsuperscript{(28)} och genom att gradvis skruva åt med den
\VUgras{stora skruven} \textsuperscript{(29)} pressar man ut saften som rinner i
ett \VUgras{litet kar} \textsuperscript{(32)} och man får ännu \VUgras{vin}
\textsuperscript{(31)}.

%%K t08-tit-8 sub
\VUsaut{6.0mm}
\VUpk{10.2pt}{40}{Öl och cider.}
\VUsaut{3.4mm}

%%K t08-c2-06-1 p
6. --- Mot \VUgras{källarens} \textsuperscript{(27)} mur klättrar en gren
\VUgras{humle} \textsuperscript{(30)}. Där är den bara en prydnadsväxt, men i
somliga länder, där vinstocken är mycket sällsynt, odlas humlen för att göra
\VUgras{öl}.

%%K t08-c2-07-1 p
7. --- Det lilla \VUgras{äppelträdet} \textsuperscript{(33)} är lastat med äpplen
som, när de är mogna, skall ge en utmärkt \VUgras{cider}. I sin \VUgras{bur}
\textsuperscript{(34)} ser \VUgras{kanariefågeln} \textsuperscript{(35)} och
\VUgras{steglitsen} \textsuperscript{(36)} med avundsjuka ögon på sparvarna som
fritt leker.

%%K t08-c2-08-1 p
8. --- Så långt ögat når, står på kullarnas sluttning \VUgras{vinstörar}
\textsuperscript{(40)} i rad som bär de praktfulla \VUgras{stammarna}
\textsuperscript{(39)}, prydda med tunga \VUgras{klasar}.

%%K t08-c2-08-2 p
Hela året har \VUgras{vinodlaren} \textsuperscript{(42)} arbetat för att sköta
sin \VUgras{vingård} \textsuperscript{(38)} och nu belönas han för sin möda med
en vacker skörd. \VUgras{Vinskörderskorna} \textsuperscript{(41)} fyller karen
som står på kärran som dras av \VUgras{mulåsnor} \textsuperscript{(43)}, och alla
är glada.

%%K t08-tit-9 sub
\VUsaut{5.0mm}
\VUpk{10.2pt}{40}{Fisket.}
\VUsaut{3.4mm}

%%K t08-c2-09-1 p
9. --- Så är det också med \VUgras{metarna} \textsuperscript{(44)} som sedan
morgonen kom och slog sig ner vid vattnets strand med sina \VUgras{metspön}
\textsuperscript{(45)}.

%%K t08-c2-09-2 p
\VUgras{Fiskarna} \textsuperscript{(47)} nappar på \VUgras{kroken} så snart de
doppar sin \VUgras{lina} \textsuperscript{(46)} i vattnet, och knappt hann de
förnya \VUgras{betet}, förrän ett nytt offer sprattlar vid linans ände.

%%K t08-c2-09-3 p
Ändå fångar \VUgras{nätfiskaren} \textsuperscript{(48)} som, från sin
\VUgras{båt} \textsuperscript{(49)}, kastar \VUgras{kastnätet} mitt i
\VUgras{floden} \textsuperscript{(50)} fler fiskar.

%%K t08-c2-10-1 p
10. --- Hösten är de vackra promenadernas årstid : till fots, till \VUgras{häst}
\textsuperscript{(52)}, på \VUgras{cykel} \textsuperscript{(53)}, med
\VUgras{bil} \textsuperscript{(54)}. \VUgras{Vägarna} \textsuperscript{(51)} är
rena och man njuter av de sista vackra dagarna, ty här är det att
\VUgras{svalorna} \textsuperscript{(56)} redan samlas på taken för att med alla
vingar flyga till varmare länder. Det gamla \VUgras{valnötsträdets}
\textsuperscript{(57)} lövverk gulnar, \VUgras{nötterna} faller och de höga
\VUgras{träden} som står grupperade som en \VUgras{dunge}
\textsuperscript{(58)} på \VUgras{höjden} \textsuperscript{(59)} skall snart
tappa sina löv.

%%K t08-c2-11-1 p
11. --- \VUgras{Solen} \textsuperscript{(60)} går upp senare, dagarna minskar, en
ganska bitande kyla börjar kännas, på \VUgras{morgonen}, och här är det att redan
flockar av \VUgras{morkullor} \textsuperscript{(61)} lämnar våra ogästvänliga
trakter.
\VUsaut{6.0mm}
\VUfilet{22mm}
